% Part: proof-theory
% Chapter: cut-elimination
% Section: ce-largest

\documentclass[../../../include/open-logic-section]{subfiles}

\begin{document}

\olfileid{pt}{cut}{inv}

\olsection{સર્વોચ્ચ ક્રમના કટ દૂર કરવાં}

\olref[seq][inv]{lem:G3c-invert}ની સાબિતી બતાવે છે કે જો
\Log{G3c} કોઈ તાર્કિક નિયમના નિષ્કર્ષને !!{prove}s
(\RightR{\lexists} અને \LeftR{\lforall} સિવાય), તો તે
સાબિતીની !!{height} વધાર્યા વિના તે નિયમનાં દરેક આધારવિધાનને
પણ !!{prove}s. આથી વધુ બતાવી શકાય છે: આ વાત
$\Log{G3c} + \CutCS$ માટે પણ સાચી છે.

અગાઉની જેમ, \CutCS{} અનુમાનનો \emph{કટ-ક્રમ} તેના
કટ !!{formula}ની ઊંડાઈ ગણીએ છીએ.

\begin{lem}\ollabel{lem:inv-G3c-cut}
જો $R$ એ \RightR{\lexists} કે \LeftR{\lforall}
સિવાયનો તાર્કિક નિયમ હોય અને
$\Log{G3c} + \CutCS$ !!a{proof}~$\pi$ વડે તેના
નિષ્કર્ષને !!{prove}s, તો તે તેના દરેક આધારવિધાનને પણ
!!{proof}~$\pi'$ વડે !!{prove}s (નિયમને બે
આધારવિધાનો હોય તો !!{proof}s~$\pi'$ અને~$\pi''$ વડે).
વધુમાં, (a)~$\pheight{\pi'}\le\pheight{\pi}$
(અને $\pheight{\pi''}\le\pheight{\pi}$) તથા
(b)~$\pi'$માં (અને $\pi''$માં) \CutCS{} અનુમાનોની
સંખ્યા તેમજ તેમના ક્રમો~$\pi$માં જેટલાં જ હોય છે.
\end{lem}

\begin{proof}
\cref{pt:seq:inv:lem:G3c-invert,pt:seq:inv:lem:invert-quant}ની
સાબિતીઓ તપાસવાથી આ મળે છે. ત્યાં
સાબિતી~$\pheight{\pi}$ પર આગમનથી હતી. આધારકિસ્સામાં
એક સ્વયંસિદ્ધને બીજા સ્વયંસિદ્ધથી બદલ્યું હતું.
આથી (a) અને~(b) સ્પષ્ટ રીતે સાચાં છે.
જો~$\pi$નો છેલ્લો નિયમ~$R$ હોય, તો તેની તત્કાળ
ઉપ-!!{proof}s જ જોઈતી !!{proof}s~$\pi_i$ છે,
અને (a) તથા~(b) બંને સાચાં છે. બીજા બધા
કિસ્સાઓમાં~$\pi$ની તત્કાળ ઉપ-!!{proof}s,
એટલે છેલ્લા અનુમાન~$R$નાં આધારવિધાનો તરફ દોરી જતી
ઉપસાબિતીઓ, પર આગમન-પૂર્વધારણા લગાવી
અને પરિણામો પર એ જ નિયમ~$R$ ફરી લગાવ્યો.
નવી સાબિતીની તત્કાળ ઉપ-!!{proof}s માટે (a) અને~(b)
સાચાં છે, તેથી સમગ્ર સાબિતી માટે પણ સાચાં છે.
છેલ્લું અનુમાન~\CutCS{} હોય તે કિસ્સો ચકાસવાનો રહે છે.
અહીં ફરી માત્ર એક ઉદાહરણ લઈએ: માનો કે
$\pi$ એ \LeftR{\land}ના નિષ્કર્ષ
$!B \land !C, \Gamma \Sequent \Delta$ની
!!a{proof} છે અને તેનો અંત~\CutCS{}થી થાય છે:
\[
\AxiomC{}
\RightLabel{$\pi_1$}
\Deduce$!B \land !C, \Gamma \fCenter \Delta, !A$
\AxiomC{}
\RightLabel{$\pi_2$}
\Deduce$!A, !B \land !C, \Gamma \fCenter \Delta$
\RightLabel{\CutCS}
\BinaryInf$!B \land !C, \Gamma \fCenter \Delta$
\DisplayProof
\]
આગમન-પૂર્વધારણા $\pi_1$ અને~$\pi_2$ પર લાગુ પડે છે
(તેઓ પણ $\LeftR{\land}$ના નિષ્કર્ષનાં દૃષ્ટાંતોની
!!{proof}s છે). તે અનુક્રમે
$!B, !C, \Gamma \fCenter \Delta, !A$ અને
$!A, !B, !C, \Gamma \fCenter \Delta$ની
!!{proof}s~$\pi_1'$ અને~$\pi_2'$ આપે છે.
તેમને \CutCS{} વડે જોડતાં~$\pi'$ મળે છે:
\[
\AxiomC{}
\RightLabel{$\pi_1'$}
\Deduce$!B, !C, \Gamma \fCenter \Delta, !A$
\AxiomC{}
\RightLabel{$\pi_2'$}
\Deduce$!A, !B, !C, \Gamma \fCenter \Delta$
\RightLabel{\CutCS}
\BinaryInf$!B, !C, \Gamma \fCenter \Delta$
\DisplayProof
\]
અહીં (a) અને~(b) સ્પષ્ટ રીતે સચવાય છે.
\end{proof}

ધ્યાન આપો કે \CutCS{}ને બદલે \Cut{} વાપર્યું હોત
તો આ રીત ચાલત નહીં. ત્યારે છેલ્લી !!{proof}ના
નિષ્કર્ષના પૂર્વાંગમાં $!B, !C, \Gamma$ની બે નકલો
અને ઉત્તરાંગમાં $\Delta$ની બે નકલો હોત.
સંકોચનની સ્વીકાર્યતા લેવી પડત, પરંતુ
\olref[seq][inv]{prop:G3c-cont-adm}ની સાબિતી
ઉલટાવી શકવાની ઉપપ્રમેયનો ઉપયોગ કરે છે.
એટલે દલીલ ચક્રીય બની જાય.

\cref{pt:cut:inv:lem:inv-G3c-cut} વડે કટ-નિરાકરણની
વૈકલ્પિક સાબિતી આપી શકાય. તેમાં \CutCS{}
અનુમાનોના સૌથી ઉપરના દૃષ્ટાંતોને એક પછી એક
દૂર કરવાને બદલે, સર્વોચ્ચ ક્રમનાં \CutCS{}
અનુમાનોને દૂર કરીએ છીએ. એટલે,
$\Log{G3c} + \CutCS$માં !!a{proof}~$\pi$ લઈને
તેમાં ક્યાંય પણ રહેલો કટ દૂર કરીએ
(કદાચ તેના સ્થાને નીચા ક્રમના કટ આવે)
અને કટ ન રહે ત્યાં સુધી આ પ્રક્રિયા ચાલુ રાખીએ.
માત્ર એટલી શરત છે કે જે કટને લઈએ તેની ઉપર
સમાન કે વધુ ક્રમનો કોઈ કટ ન હોય.

\begin{lem}\ollabel{lem:max-cut-red-G3c}
જો $\pi$ એ $n$ ક્રમના~\CutCS{} અનુમાનથી પૂરી થતી
$\Log{G3c}+\CutCS$ની !!a{proof} હોય, અને તે ઉપરાંત
માત્ર~$\Log{G3c}$ના નિયમો તથા $<n$ ક્રમનાં
\CutCS{} અનુમાનો વાપરતી હોય, તો એ જ અંતિમ
સિક્વન્ટની એવી સાબિતી~$\pi'$ પણ
$\Log{G3c} + \CutCS$માં છે જેમાં દરેક
\CutCS{} અનુમાનનો ક્રમ~$<n$ છે.
\end{lem}

\begin{proof}
!!{proof}~$\pi$નો અંત આ રીતે થાય છે:
\[
\AxiomC{}
\RightLabel{$\pi_1$}
\Deduce$\Gamma \fCenter \Delta, !A$
\AxiomC{}
\RightLabel{$\pi_2$}
\Deduce$!A, \Gamma \fCenter \Delta$
\RightLabel{\CutCS}
\BinaryInf$\Gamma \fCenter \Delta$
\DisplayProof
\]
હવે~$!A$ના સ્વરૂપ પ્રમાણે કિસ્સા જુદા પાડીએ.

માનો કે~$!A$ પરમાણ્વીય છે. શરત મુજબ
~$\pi_1$ અને $\pi_2$માંનાં બધાં \CutCS{} અનુમાનોના
ક્રમ $<\depth{!A} = 0$ હોવા જોઈએ;
આથી~$\pi_1$ કે~$\pi_2$માં \CutCS{} અનુમાન હોઈ જ શકે નહીં.
$\pheight{\pi_2}$ પર આગમનથી બતાવીએ કે
$\Gamma \Sequent \Delta$ની કટ-મુક્ત સાબિતી છે.
જો $\pheight{\pi_2} = 0$, તો
$!A,\Gamma \Sequent \Delta$ સ્વયંસિદ્ધ છે.
જો~$!A$ મુખ્ય ન હોય, તો કોઈ પરમાણ્વીય
$!C$ એ~$\Gamma$ તથા $\Delta$ બંનેનું
!!a{element} છે, એટલે
$\Gamma \Sequent \Delta$ પોતે સ્વયંસિદ્ધ છે.
જો~$!A$ મુખ્ય હોય, તો
$\Delta = \Delta', !A$.
\sourcecorrection{OLMD-196}{મૂળમાં $\Delta = \Delta, !A$ છે; અહીં બાકીના ઉત્તરાંગ માટે $\Delta'$ જરૂરી છે.}
આ કિસ્સામાં~$\pi_1$નો અંત
$\Gamma \Sequent \Delta', !A, !A$થી થાય છે.
\olref[seq][inv]{prop:G3c-cont-adm}
(સંકોચનની સ્વીકાર્યતા) લગાવીને
$\Gamma \Sequent \Delta', !A$ની
કટ-મુક્ત સાબિતી મળે છે.
જો $\pheight{\pi_2} > 0$, તો તેનો અંત
કોઈ નિયમ~$R$થી થાય છે. તે નિયમનું એક
આધારવિધાન લો: તેનું સ્વરૂપ
$!A, \Gamma',\Pi \Sequent \Delta',\Lambda$
છે, જ્યાં $\Pi$ અને $\Lambda$ એ~$R$નાં
સક્રિય !!{formula}s છે.
\olref[seq][adm]{prop:weak-G3c-adm}ને
~$\pi_1$ અને~$\pi_2$નાં યોગ્ય તત્કાળ
ઉપ-સાબિતીઓ પર લગાવતાં
$\Gamma,\Gamma',\Pi \Sequent
\Delta,\Delta',\Lambda,!A$ અને
$!A,\Gamma,\Gamma',\Pi \Sequent
\Delta,\Delta',\Lambda$ની !!{proof}s મળે છે.
\sourcecorrection{OLMD-197}{મૂળ લખાણ $\pi_2$ પર સીધું શિથિલીકરણ કહે છે; નિયમના આધારવિધાન સુધી પહોંચવા તેની યોગ્ય તત્કાળ ઉપ-સાબિતી જોઈએ.}
તેના પર આગમન-પૂર્વધારણા લાગુ કરતાં
~$R$ના દરેક આધારવિધાન માટે
$\Gamma,\Gamma',\Pi \Sequent
\Delta,\Delta',\Lambda$ની !!a{proof} મળે છે.
પછી~$R$ લગાવતાં
$\Gamma,\Gamma \Sequent \Delta,\Delta$ મળે છે;
\olref[seq][inv]{prop:G3c-cont-adm} વડે
$\Gamma \Sequent \Delta$ની કટ-મુક્ત
સાબિતી મળે છે.

જો~$!A$ પરમાણ્વીય ન હોય, તો તેના
!!{main
operator} પ્રમાણે~$!A$ના કિસ્સા જુદા પાડીએ.
દરેક કિસ્સામાં \olref{lem:inv-G3c-cut}
લગાવીને, $!A$ મુખ્ય સૂત્ર હોય તેવા
ડાબા અને જમણા નિયમોના આધારવિધાનોની
!!{proof}s મેળવીએ. પછી
\olref[top]{lem:cut-adm-G3c}ની સાબિતીના
કિસ્સા~E મુજબ આગળ વધીએ.
\begin{enumerate}
  \item $!A \ident \lnot !B$. !!{proof}s $\pi_1$ અને $\pi_2$નો
  અંત $\Gamma \Sequent \Delta, \lnot !B$ અને
  $\lnot !B, \Gamma \Sequent \Delta$થી થાય છે.
  \olref{lem:inv-G3c-cut} મુજબ
  $!B, \Gamma \Sequent \Delta$ અને
  $\Gamma \Sequent \Delta, !B$ની
  !!{proof}s $\pi_1'$ અને $\pi_2'$ છે.
  !!{proof}~$\pi'$ આ છે:
  \[
  \AxiomC{}
  \RightLabel{$\pi_2'$}
  \Deduce$\Gamma \fCenter \Delta, !B$
  \AxiomC{}
  \RightLabel{$\pi_1'$}
  \Deduce$!B, \Gamma \fCenter \Delta$
  \RightLabel{\CutCS}
  \BinaryInf$\Gamma \fCenter \Delta$
  \DisplayProof
  \]
  \item $!A \ident (!B \lor !C)$. !!{proof}s $\pi_1$ અને $\pi_2$નો
  અંત $\Gamma \Sequent \Delta, !B \lor !C$ અને
  $!B \lor !C, \Gamma \Sequent \Delta$થી થાય છે.
  \olref{lem:inv-G3c-cut} વડે
  ~$\pi_1$ પર \RightR{\lor}નું ઉલટાવવું
  લગાવતાં $\Gamma \Sequent \Delta, !B, !C$ની
  !!a{proof}~$\pi_1'$ મળે છે.
  \olref{lem:inv-G3c-cut} વડે~$\pi_2$ પર
  \LeftR{\lor}નું ઉલટાવવું લગાવતાં
  $!B, \Gamma \Sequent \Delta$ની
  !!a{proof}~$\pi_2'$ અને
  $!C, \Gamma \Sequent \Delta$ની
  !!a{proof}~$\pi_2''$ મળે છે.
  $\pi_2''$ પર
  \olref[seq][adm]{prop:weak-G3c-adm}
  લગાવતાં $!C, \Gamma \Sequent \Delta, !B$ની
  !!a{proof}~$\pi_2'''$ પણ મળે છે.
  !!{proof}~$\pi'$ આ છે:
  \[
  \AxiomC{}
  \RightLabel{$\pi_1'$}
  \Deduce$\Gamma \fCenter \Delta, !B, !C$
  \AxiomC{}
  \RightLabel{$\pi_2'''$}
  \Deduce$!C, \Gamma \fCenter \Delta, !B$
  \RightLabel{\CutCS}
  \BinaryInf$\Gamma \fCenter \Delta, !B$
  \AxiomC{}
  \RightLabel{$\pi_2'$}
  \Deduce$!B, \Gamma \fCenter \Delta$
  \RightLabel{\CutCS}
  \BinaryInf$\Gamma \fCenter \Delta$
  \DisplayProof
  \]
  \item $!A \ident (!B \land !C)$. કસરત.
  \item $!A \ident (!B \lif !C)$. !!{proof}s $\pi_1$ અને $\pi_2$નો
  અંત $\Gamma \Sequent \Delta, !B \lif !C$ અને
  $!B \lif !C, \Gamma \Sequent \Delta$થી થાય છે.
  \olref{lem:inv-G3c-cut} વડે
  ~$\pi_1$ પર \RightR{\lif}નું ઉલટાવવું
  લગાવતાં $!B, \Gamma \Sequent \Delta, !C$ની
  !!a{proof}~$\pi_1'$ મળે છે.
  \olref{lem:inv-G3c-cut} વડે~$\pi_2$ પર
  \LeftR{\lif}નું ઉલટાવવું લગાવતાં
  $\Gamma \Sequent \Delta, !B$ની
  !!a{proof}~$\pi_2'$ અને
  $!C, \Gamma \Sequent \Delta$ની
  !!a{proof}~$\pi_2''$ મળે છે.
  $\pi_2''$ પર
  \olref[seq][adm]{prop:weak-G3c-adm}
  લગાવતાં $!C, !B, \Gamma \Sequent \Delta$ની
  !!a{proof}~$\pi_2'''$ પણ મળે છે.
  !!{proof}~$\pi'$ આ છે:
  \[
  \AxiomC{}
  \RightLabel{$\pi_2'$}
  \Deduce$\Gamma \fCenter \Delta, !B$
  \AxiomC{}
  \RightLabel{$\pi_1'$}
  \Deduce$!B, \Gamma \fCenter \Delta, !C$
  \AxiomC{}
  \RightLabel{$\pi_2'''$}
  \Deduce$!C, !B, \Gamma \fCenter \Delta$
  \RightLabel{\CutCS}
  \BinaryInf$!B, \Gamma \fCenter \Delta$
  \RightLabel{\CutCS}
  \BinaryInf$\Gamma \fCenter \Delta$
  \DisplayProof
  \]
  \item $!A \ident \lforall[x][!B(x)]$. !!{proof}s $\pi_1$ અને
  $\pi_2$નો અંત $\Gamma \Sequent \Delta, \lforall[x][!B(x)]$
  અને $\lforall[x][!B(x)], \Gamma \Sequent \Delta$થી થાય છે.
  \olref{lem:inv-G3c-cut} મુજબ
  $\Gamma \Sequent \Delta, !B(c)$ની
  !!a{proof}~$\pi_1'$ છે.

  હવે !!{proof}~$\pi_2$ લો. તેનો અંત
  $\lforall[x][!B(x)], \Gamma \Sequent \Delta$થી
  થાય છે. $\pi_2$ના અંતિમ સિક્વન્ટથી ઉપરની તરફ
  જતાં, સિક્વન્ટની ડાબી બાજુ રહેલા
  $\lforall[x][!B(x)]$ના એવા દરેક દૃષ્ટાંતને
  ચિહ્નિત કરો જે $\pi_2$ના અંતિમ સિક્વન્ટમાં
  એ જ $\lforall[x][!B(x)]$ના દૃષ્ટાંત સુધી પહોંચે છે. એટલે:
  (a)~$\pi_2$ના અંતિમ સિક્વન્ટમાં
  $\lforall[x][!B(x)]$ ચિહ્નિત કરો.
  (b)~નિયમના નિષ્કર્ષના સંદર્ભમાં
  $\lforall[x][!B(x)]$
  ચિહ્નિત હોય તો તેના આધારવિધાનમાં અનુરૂપ
  દૃષ્ટાંત ચિહ્નિત કરો. (c)~\LeftR{\lforall}
  નિયમના નિષ્કર્ષમાં $\lforall[x][!B(x)]$
  મુખ્ય અને ચિહ્નિત હોય, તો આધારવિધાનમાં
  $\lforall[x][!B(x)]$ તથા $!B(t)$ના
  અનુરૂપ સક્રિય દૃષ્ટાંતો ચિહ્નિત કરો.
  
  હવે $\lforall[x][!B(x)]$ના આ બધા
  ચિહ્નિત દૃષ્ટાંતો લો. \LeftR{\lforall}
  સિવાયના કોઈ નિયમમાં તેઓ સક્રિય નથી
  (ચિહ્નિત થનારા સક્રિય !!{formula}s
  માત્ર \LeftR{\lforall}નાં છે).
  $\pi_2$માંથી $\lforall[x][!B(x)]$ના
  બધા ચિહ્નિત દૃષ્ટાંતો કાઢી નાખો.
  જો મળેલું લખાણ સાચી !!{proof} હોય,
  તો $\lforall[x][!B(x)]$ના ચિહ્નિત
  દૃષ્ટાંતો કદી સક્રિય નહોતા;
  તેઓ બધા સીધા સ્વયંસિદ્ધોમાંથી આવ્યા હતા.
  !!{proof}નો અંત~$\Gamma \Sequent \Delta$થી
  થાય છે, તેથી~$\pi$ને તેનાથી બદલી શકાય.
  આમ સર્વોચ્ચ ક્રમનો એક કટ દૂર થયો.

  અન્યથા, મળેલું લખાણ સાચી સાબિતી નથી,
  કારણ કે કેટલીક \LeftR{\lforall} નિષ્પત્તિઓમાં
  સક્રિય રહેલાં $\lforall[x][!B(x)]$
  !!{formula}s કાઢી નાખ્યાં છે.
  તેથી તે નિષ્પત્તિઓ આ સ્વરૂપની
  ``નિષ્પત્તિઓ'' બની છે:
  \[
  \AxiomC{}
  \RightLabel{$\pi_t$}
  \Deduce$!B(t), \Pi \fCenter \Lambda$
  \RightLabel{\LeftR{\lforall}}
  \UnaryInf$\Pi \fCenter \Lambda$
  \DisplayProof
  \]
  આવા સૌથી ઉપરના દરેક ``અનુમાન''ને
  આથી બદલો:
  \[
  \AxiomC{}
  \RightLabel{$\Subst{\pi_1'}{t}{c}[\Pi \Sequent \Lambda]$}
  \Deduce$\Gamma, \Pi \fCenter \Delta, \Lambda, !B(t)$
  \AxiomC{}
  \RightLabel{$\pi_t[\Gamma \Sequent \Delta]$}
  \Deduce$!B(t), \Gamma, \Pi \fCenter \Delta, \Lambda$
  \RightLabel{\CutCS}
  \BinaryInf$\Gamma, \Pi \fCenter \Delta, \Lambda$
  \DisplayProof
  \]
  અને તેની નીચેના દરેક સિક્વન્ટની ડાબે
  $\Gamma$ તથા જમણે $\Delta$ ઉમેરો.
  આધારવિધાનમાં ચિહ્નિત $!B(t)$ ધરાવતાં
  બધા \LeftR{\lforall} ``અનુમાનો''ને
  કોઈ~$!B(t)$ પરના \CutCS{} અનુમાનથી
  બદલી ન દેવાય ત્યાં સુધી આ દોહરાવો.
  મળેલી !!{proof}નો અંતિમ સિક્વન્ટ
  (અને હવે તે સાચી !!{proof} છે)
  $\Gamma, \dots, \Gamma \Sequent
  \Delta, \dots, \Delta$ સ્વરૂપનો છે.
  સંકોચનની સ્વીકાર્યતા લગાવી તેને
  $\Gamma \Sequent \Delta$ની !!a{proof}માં
  ફેરવો અને~$\pi$ને તેનાથી બદલો.
  $\lforall[x][!B(x)]$ પરના એક કટને
  $!B(t)$ પરના (કદાચ ઘણા) કટથી બદલ્યો છે,
  પણ તે બધા નીચા ક્રમના છે.
  $\pi_1$ કે $\pi_2$માં પહેલેથી રહેલા
  બધા કટ પણ રહે છે, અને તે પણ
  નીચા ક્રમના જ છે.
\end{enumerate}
\end{proof}

\begin{prob}
\olref{lem:max-cut-red-G3c}ની સાબિતીમાં
$!A \ident (!B \land !C)$નો કિસ્સો ચકાસો.
\sourcecorrection{OLMD-198}{મૂળ કસરત ઉલટાવવાના ઉપપ્રમેયનો સંદર્ભ આપે છે; અહીં ચર્ચાતો કિસ્સો મહત્તમ-ક્રમના કટ ઘટાડવાના ઉપપ્રમેયનો છે.}
અર્થાત્, જો અંતમાં
\[
\AxiomC{}
\RightLabel{$\pi_1$}
\Deduce$\Gamma \fCenter \Delta, !B \land !C$
\AxiomC{}
\RightLabel{$\pi_2$}
\Deduce$!B \land !C, \Gamma \fCenter \Delta$
\RightLabel{\CutCS}
\BinaryInf$\Gamma \fCenter \Delta$
\DisplayProof
\]
ધરાવતી !!a{proof}~$\pi$ હોય અને
$\pi_1$ તથા $\pi_2$માં
$<\depth{!B \land !C}$ ક્રમના કટ જ હોય,
તો $\Gamma \Sequent \Delta$ની એવી
!!a{proof}~$\pi'$ છે જેમાં પણ
$<\depth{!B \land !C}$ ક્રમના કટ જ છે,
એ બતાવો.
\end{prob}

% \begin{prob}
% Verify the case where $!A \ident \lexists[x][!B(x)]$ in the proof of
% \olref{lem:inv-G3c-cut}.
% \end{prob}


\olref{lem:max-cut-red-G3c} બતાવે છે કે
!!a{proof}માં સર્વોચ્ચ ક્રમના \CutCS{}
અનુમાનને નીચા ક્રમનાં \CutCS{} અનુમાનોથી
બદલી શકાય.
\sourcecorrection{OLMD-199}{મૂળ અનુચ્છેદ અગાઉના બીજા ઉપપ્રમેય તરફ નિર્દેશ કરે છે; અહીં જરૂરી પરિણામ આ વિભાગનું મહત્તમ-ક્રમના કટ ઘટાડવાનું ઉપપ્રમેય છે.}
આ ઉપપ્રમેય ફરી લગાવી,
સર્વોચ્ચ ક્રમના~\CutCS{} પર સમાપ્ત
થતી સૌથી ઉપરની ઉપ-!!{proof}s પર
ક્રમશઃ પ્રક્રિયા કરીને, કોઈ પણ
!!a{proof}માંથી ગમે તેટલા \CutCS{}
અનુમાનો દૂર કરી શકાય છે.

\begin{thm}\ollabel{thm:cut-elim-G3c-largest}
$\Log{G3c} + \CutCS$ની કોઈ પણ
!!{proof}~$\pi$ને~\Log{G3c}ની
સાબિતીમાં ફેરવી શકાય છે.
\end{thm}

\begin{proof}
  પહેલાં બતાવીએ કે~$\pi$માંના સર્વોચ્ચ
  ક્રમના બધા કટ દૂર કરી શકાય છે.
  માનો કે~$d$ એ~$\pi$માંના કટનો
  મહત્તમ ક્રમ છે અને~$n$ એ~$d$
  ક્રમના કટની સંખ્યા છે. સાબિતી
  ~$n$ પર આગમનથી છે.
  જો $n=0$, તો કશું કરવાનું નથી.
  જો $n>0$, તો~$d$ ક્રમનો
  સૌથી ઉપરનો કટ પસંદ કરો અને
  તેમાં પૂરી થતી ઉપ-!!{proof}ને
  $\pi_1$ કહો.
  \olref{lem:max-cut-red-G3c} મુજબ
  એ જ અંતિમ સિક્વન્ટની એવી
  !!a{proof}~$\pi_1'$ છે જેમાં
  બધા કટનો ક્રમ $<d$ છે.
  $\pi$માં $\pi_1$ને $\pi_1'$થી
  બદલો. મળેલી !!a{proof}માં
  $d$ ક્રમના $n-1$ કટ રહે છે,
  તેથી આગમન-પૂર્વધારણા લાગુ પડે છે.

  હવે~$d$ પર આગમનથી પરિણામ મળે છે.
  જો $d=0$, તો બધા કટ પરમાણ્વીય છે
  અને અગાઉના પરિણામથી તે બધા
  દૂર કરી શકાય છે. જો $d>0$,
  તો અગાઉનું પરિણામ એવી સાબિતી આપે
  છે જેમાં $d$ ક્રમના બધા કટ
  $<d$ ક્રમના કટથી બદલાયા છે.
  કારણ કે~$d$ એ~$\pi$માંના કટનો
  મહત્તમ ક્રમ હતો, મળેલી સાબિતીમાં
  કટનો મહત્તમ ક્રમ $<d$ છે,
  અને આગમન-પૂર્વધારણા લાગુ પડે છે.
\end{proof}

\end{document}
